\section{Introduction and Motivation}

Quantum algorithms are widely expected to offer an eventual advantage over classical methods for describing quantum many-body systems, especially in electronic-structure problems where the Hilbert space grows exponentially with system size. Quantum phase estimation (QPE) is one of the representative examples because, in principle, it can recover eigenvalues to high precision when provided with a sufficiently good initial state \cite{abrams1999qpe}. In practice, however, QPE requires deep coherent time evolution and fault-tolerant resources that remain beyond the reach of current hardware for most chemically relevant systems.

These limitations motivated the development of hybrid algorithms such as the variational quantum eigensolver (VQE) \cite{peruzzo2014vqe}. VQE replaces long coherent evolution with repeated preparation and measurement of a parameterized trial state, while a classical optimizer updates the circuit parameters to minimize the measured energy. This hybrid structure is much more compatible with near-term hardware, but it introduces a different bottleneck: the variational optimization itself can be costly, sensitive to initialization, and difficult to scale when the ansatz contains many parameters.

\subsection{Why It Matters}

As an alternative, subspace methods have also been developed for molecular ground-state estimation. Instead of optimizing a single wavefunction, these methods prepare a small set of basis states and solve a projected eigenvalue problem within the span of those states. In this setting, the quality of the chosen basis becomes the central issue. A better basis can improve the accuracy of the projected ground-state estimate, while a lower-cost basis can reduce circuit depth, gate count, and measurement overhead.

This motivates the present study. Among the many basis-construction strategies proposed in the literature, Krylov states generated by Hamiltonian evolution \cite{cortes2022qksd} and UCCSD states generated by chemically motivated single and double excitations \cite{anand2022ucc} are two representative state-of-the-art directions. Krylov uses the Hamiltonian itself to generate candidate directions in Hilbert space, whereas UCCSD uses an excitation structure inherited from coupled-cluster theory. Although many variations of both approaches exist, these two families serve here as representative Hamiltonian-guided and chemistry-guided basis constructions.

\subsection{Problem Statement}

There are several distinct ways to evaluate a quantum state-preparation strategy for ground-state estimation. In a QPE-type setting, a prepared state is used as the initial state for phase estimation, so its quality is judged mainly by its overlap with the target eigenstate. In a VQE-type setting, the state or ansatz is used as the trial form in a variational optimization, so its quality is judged by the minimum energy reachable after parameter updates. In a subspace method, by contrast, a family of prepared states is treated as a basis set, and the quality of that family is judged by the projected eigensolver built from their span.

This work adopts the third viewpoint. The main object of interest in this paper is not just one trial state, but the basis generator itself. Evaluating Krylov and UCCSD through the subspace they generate makes it possible to compare their collective expressivity, which can be richer than what is visible from a single state alone. This choice is especially important for Krylov states: a single Krylov state of the form
\begin{equation}
\ket{\phi(t)} = e^{-it\hat{H}}\ket{\phi_{\mathrm{HF}}}
\end{equation}
has a constant expectation value of $\hat{H}$ under the same Hamiltonian, so it is not meaningful to assess Krylov as a one-state variational ansatz. Krylov is naturally a set-based construction, and subspace projection is therefore the appropriate framework for comparing it with UCCSD as a basis-state generator. The concrete problem studied in this paper is molecular ground-state energy estimation. Starting from the same Hartree-Fock reference and using the same target subspace size $K$, this study asks which basis family yields the lower projected ground-state energy for a given state-preparation cost, with final projected energies compared against exact full configuration interaction (FCI) reference data on benchmark molecules.
